\documentclass{report}

%%%%%%%%%%%%%%%%%%%%%%%%%%%%%%%%%
% PACKAGE IMPORTS
%%%%%%%%%%%%%%%%%%%%%%%%%%%%%%%%%
\usepackage[dvipsnames]{xcolor}


\usepackage{anyfontsize}
\usepackage{braket}


\usepackage[tmargin=2cm,rmargin=1in,lmargin=1in,margin=0.85in,bmargin=2cm,footskip=.2in]{geometry}
\usepackage{amsmath,amsfonts,amsthm,amssymb,mathtools}
%\usepackage[varbb]{newpxmath}
\usepackage{xfrac}
\usepackage[makeroom]{cancel}
\usepackage{mathtools}
\usepackage{bookmark}
\usepackage{enumitem}
\usepackage{hyperref,theoremref}
\hypersetup{
	pdftitle={Assignment},
	colorlinks=true, linkcolor=doc!90,
	bookmarksnumbered=true,
	bookmarksopen=true
}
\usepackage[most,many,breakable]{tcolorbox}
\usepackage{xcolor}
\usepackage{varwidth}
\usepackage{varwidth}
\usepackage{etoolbox}
%\usepackage{authblk}
\usepackage{nameref}
\usepackage{multicol,array}
\usepackage{tikz-cd}
\usepackage[ruled,vlined,linesnumbered]{algorithm2e}
\usepackage{comment} % enables the use of multi-line comments (\ifx \fi) 
\usepackage{import}
\usepackage{xifthen}
\usepackage{pdfpages}
\usepackage{transparent}

\usetikzlibrary{shapes.geometric}
\usetikzlibrary{shapes}

\usetikzlibrary{calc}

\newcommand\mycommfont[1]{\footnotesize\ttfamily\textcolor{blue}{#1}}
\SetCommentSty{mycommfont}
\newcommand{\incfig}[1]{%
    \def\svgwidth{\columnwidth}
    \import{./figures/}{#1.pdf_tex}
}

\usepackage{tikzsymbols}
\renewcommand\qedsymbol{\sus}


%\usepackage{import}
%\usepackage{xifthen}
%\usepackage{pdfpages}
%\usepackage{transparent}


%%%%%%%%%%%%%%%%%%%%%%%%%%%%%%
% SELF MADE COLORS
%%%%%%%%%%%%%%%%%%%%%%%%%%%%%%



\definecolor{myg}{RGB}{56, 140, 70}
\definecolor{myb}{RGB}{45, 111, 177}
\definecolor{myr}{RGB}{199, 68, 64}
\definecolor{mytheorembg}{HTML}{F2F2F9}
\definecolor{mytheoremfr}{HTML}{00007B}
\definecolor{mylenmabg}{HTML}{FFFAF8}
\definecolor{mylenmafr}{HTML}{983b0f}
\definecolor{mypropbg}{HTML}{f2fbfc}
\definecolor{mypropfr}{HTML}{191971}
\definecolor{myexamplebg}{HTML}{F2FBF8}
\definecolor{myexamplefr}{HTML}{88D6D1}
\definecolor{myexampleti}{HTML}{2A7F7F}
\definecolor{mydefinitbg}{HTML}{E5E5FF}
\definecolor{mydefinitfr}{HTML}{3F3FA3}
\definecolor{notesgreen}{RGB}{0,162,0}
\definecolor{myp}{RGB}{197, 92, 212}
\definecolor{mygr}{HTML}{2C3338}
\definecolor{myred}{RGB}{127,0,0}
\definecolor{myyellow}{RGB}{169,121,69}
\definecolor{myexercisebg}{HTML}{F2FBF8}
\definecolor{myexercisefg}{HTML}{88D6D1}


%%%%%%%%%%%%%%%%%%%%%%%%%%%%
% TCOLORBOX SETUPS
%%%%%%%%%%%%%%%%%%%%%%%%%%%%

\setlength{\parindent}{1cm}
%================================
% THEOREM BOX
%================================

\tcbuselibrary{theorems,skins,hooks}
\newtcbtheorem[number within=section]{Theorem}{Theorem}
{%
	enhanced,
	breakable,
	colback = mytheorembg,
	frame hidden,
	boxrule = 0sp,
	borderline west = {2pt}{0pt}{mytheoremfr},
	sharp corners,
	detach title,
	before upper = \tcbtitle\par\smallskip,
	coltitle = mytheoremfr,
	fonttitle = \bfseries\sffamily,
	description font = \mdseries,
	separator sign none,
	segmentation style={solid, mytheoremfr},
}
{th}

\tcbuselibrary{theorems,skins,hooks}
\newtcbtheorem[number within=chapter]{theorem}{Theorem}
{%
	enhanced,
	breakable,
	colback = mytheorembg,
	frame hidden,
	boxrule = 0sp,
	borderline west = {2pt}{0pt}{mytheoremfr},
	sharp corners,
	detach title,
	before upper = \tcbtitle\par\smallskip,
	coltitle = mytheoremfr,
	fonttitle = \bfseries\sffamily,
	description font = \mdseries,
	separator sign none,
	segmentation style={solid, mytheoremfr},
}
{th}


\tcbuselibrary{theorems,skins,hooks}
\newtcolorbox{Theoremcon}
{%
	enhanced
	,breakable
	,colback = mytheorembg
	,frame hidden
	,boxrule = 0sp
	,borderline west = {2pt}{0pt}{mytheoremfr}
	,sharp corners
	,description font = \mdseries
	,separator sign none
}

%================================
% Corollery
%================================
\tcbuselibrary{theorems,skins,hooks}
\newtcbtheorem[number within=section]{Corollary}{Corollary}
{%
	enhanced
	,breakable
	,colback = myp!10
	,frame hidden
	,boxrule = 0sp
	,borderline west = {2pt}{0pt}{myp!85!black}
	,sharp corners
	,detach title
	,before upper = \tcbtitle\par\smallskip
	,coltitle = myp!85!black
	,fonttitle = \bfseries\sffamily
	,description font = \mdseries
	,separator sign none
	,segmentation style={solid, myp!85!black}
}
{th}
\tcbuselibrary{theorems,skins,hooks}
\newtcbtheorem[number within=chapter]{corollary}{Corollary}
{%
	enhanced
	,breakable
	,colback = myp!10
	,frame hidden
	,boxrule = 0sp
	,borderline west = {2pt}{0pt}{myp!85!black}
	,sharp corners
	,detach title
	,before upper = \tcbtitle\par\smallskip
	,coltitle = myp!85!black
	,fonttitle = \bfseries\sffamily
	,description font = \mdseries
	,separator sign none
	,segmentation style={solid, myp!85!black}
}
{th}


%================================
% LENMA
%================================

\tcbuselibrary{theorems,skins,hooks}
\newtcbtheorem[number within=section]{Lenma}{Lemma}
{%
	enhanced,
	breakable,
	colback = mylenmabg,
	frame hidden,
	boxrule = 0sp,
	borderline west = {2pt}{0pt}{mylenmafr},
	sharp corners,
	detach title,
	before upper = \tcbtitle\par\smallskip,
	coltitle = mylenmafr,
	fonttitle = \bfseries\sffamily,
	description font = \mdseries,
	separator sign none,
	segmentation style={solid, mylenmafr},
}
{th}

\tcbuselibrary{theorems,skins,hooks}
\newtcbtheorem[number within=chapter]{lenma}{Lenma}
{%
	enhanced,
	breakable,
	colback = mylenmabg,
	frame hidden,
	boxrule = 0sp,
	borderline west = {2pt}{0pt}{mylenmafr},
	sharp corners,
	detach title,
	before upper = \tcbtitle\par\smallskip,
	coltitle = mylenmafr,
	fonttitle = \bfseries\sffamily,
	description font = \mdseries,
	separator sign none,
	segmentation style={solid, mylenmafr},
}
{th}


%================================
% PROPOSITION
%================================

\tcbuselibrary{theorems,skins,hooks}
\newtcbtheorem[number within=section]{Prop}{Proposition}
{%
	enhanced,
	breakable,
	colback = mypropbg,
	frame hidden,
	boxrule = 0sp,
	borderline west = {2pt}{0pt}{mypropfr},
	sharp corners,
	detach title,
	before upper = \tcbtitle\par\smallskip,
	coltitle = mypropfr,
	fonttitle = \bfseries\sffamily,
	description font = \mdseries,
	separator sign none,
	segmentation style={solid, mypropfr},
}
{th}

\tcbuselibrary{theorems,skins,hooks}
\newtcbtheorem[number within=chapter]{prop}{Proposition}
{%
	enhanced,
	breakable,
	colback = mypropbg,
	frame hidden,
	boxrule = 0sp,
	borderline west = {2pt}{0pt}{mypropfr},
	sharp corners,
	detach title,
	before upper = \tcbtitle\par\smallskip,
	coltitle = mypropfr,
	fonttitle = \bfseries\sffamily,
	description font = \mdseries,
	separator sign none,
	segmentation style={solid, mypropfr},
}
{th}


%================================
% CLAIM
%================================

\tcbuselibrary{theorems,skins,hooks}
\newtcbtheorem[number within=section]{claim}{Claim}
{%
	enhanced
	,breakable
	,colback = myg!10
	,frame hidden
	,boxrule = 0sp
	,borderline west = {2pt}{0pt}{myg}
	,sharp corners
	,detach title
	,before upper = \tcbtitle\par\smallskip
	,coltitle = myg!85!black
	,fonttitle = \bfseries\sffamily
	,description font = \mdseries
	,separator sign none
	,segmentation style={solid, myg!85!black}
}
{th}



%================================
% Exercise
%================================

\tcbuselibrary{theorems,skins,hooks}
\newtcbtheorem[number within=section]{Exercise}{Exercise}
{%
	enhanced,
	breakable,
	colback = myexercisebg,
	frame hidden,
	boxrule = 0sp,
	borderline west = {2pt}{0pt}{myexercisefg},
	sharp corners,
	detach title,
	before upper = \tcbtitle\par\smallskip,
	coltitle = myexercisefg,
	fonttitle = \bfseries\sffamily,
	description font = \mdseries,
	separator sign none,
	segmentation style={solid, myexercisefg},
}
{th}

\tcbuselibrary{theorems,skins,hooks}
\newtcbtheorem[number within=chapter]{exercise}{Exercise}
{%
	enhanced,
	breakable,
	colback = myexercisebg,
	frame hidden,
	boxrule = 0sp,
	borderline west = {2pt}{0pt}{myexercisefg},
	sharp corners,
	detach title,
	before upper = \tcbtitle\par\smallskip,
	coltitle = myexercisefg,
	fonttitle = \bfseries\sffamily,
	description font = \mdseries,
	separator sign none,
	segmentation style={solid, myexercisefg},
}
{th}

%================================
% EXAMPLE BOX
%================================

\newtcbtheorem[number within=section]{Example}{Example}
{%
	colback = myexamplebg
	,breakable
	,colframe = myexamplefr
	,coltitle = myexampleti
	,boxrule = 1pt
	,sharp corners
	,detach title
	,before upper=\tcbtitle\par\smallskip
	,fonttitle = \bfseries
	,description font = \mdseries
	,separator sign none
	,description delimiters parenthesis
}
{ex}

\newtcbtheorem[number within=chapter]{example}{Example}
{%
	colback = myexamplebg
	,breakable
	,colframe = myexamplefr
	,coltitle = myexampleti
	,boxrule = 1pt
	,sharp corners
	,detach title
	,before upper=\tcbtitle\par\smallskip
	,fonttitle = \bfseries
	,description font = \mdseries
	,separator sign none
	,description delimiters parenthesis
}
{ex}

%================================
% DEFINITION BOX
%================================

\newtcbtheorem[number within=section]{Definition}{Definition}{enhanced,
	before skip=2mm,after skip=2mm, colback=red!5,colframe=red!80!black,boxrule=0.5mm,
	attach boxed title to top left={xshift=1cm,yshift*=1mm-\tcboxedtitleheight}, varwidth boxed title*=-3cm,
	boxed title style={frame code={
					\path[fill=tcbcolback]
					([yshift=-1mm,xshift=-1mm]frame.north west)
					arc[start angle=0,end angle=180,radius=1mm]
					([yshift=-1mm,xshift=1mm]frame.north east)
					arc[start angle=180,end angle=0,radius=1mm];
					\path[left color=tcbcolback!60!black,right color=tcbcolback!60!black,
						middle color=tcbcolback!80!black]
					([xshift=-2mm]frame.north west) -- ([xshift=2mm]frame.north east)
					[rounded corners=1mm]-- ([xshift=1mm,yshift=-1mm]frame.north east)
					-- (frame.south east) -- (frame.south west)
					-- ([xshift=-1mm,yshift=-1mm]frame.north west)
					[sharp corners]-- cycle;
				},interior engine=empty,
		},
	fonttitle=\bfseries,
	title={#2},#1}{def}
\newtcbtheorem[number within=chapter]{definition}{Definition}{enhanced,
	before skip=2mm,after skip=2mm, colback=red!5,colframe=red!80!black,boxrule=0.5mm,
	attach boxed title to top left={xshift=1cm,yshift*=1mm-\tcboxedtitleheight}, varwidth boxed title*=-3cm,
	boxed title style={frame code={
					\path[fill=tcbcolback]
					([yshift=-1mm,xshift=-1mm]frame.north west)
					arc[start angle=0,end angle=180,radius=1mm]
					([yshift=-1mm,xshift=1mm]frame.north east)
					arc[start angle=180,end angle=0,radius=1mm];
					\path[left color=tcbcolback!60!black,right color=tcbcolback!60!black,
						middle color=tcbcolback!80!black]
					([xshift=-2mm]frame.north west) -- ([xshift=2mm]frame.north east)
					[rounded corners=1mm]-- ([xshift=1mm,yshift=-1mm]frame.north east)
					-- (frame.south east) -- (frame.south west)
					-- ([xshift=-1mm,yshift=-1mm]frame.north west)
					[sharp corners]-- cycle;
				},interior engine=empty,
		},
	fonttitle=\bfseries,
	title={#2},#1}{def}



%================================
% Solution BOX
%================================

\makeatletter
\newtcbtheorem{question}{Question}{enhanced,
	breakable,
	colback=white,
	colframe=myb!80!black,
	attach boxed title to top left={yshift*=-\tcboxedtitleheight},
	fonttitle=\bfseries,
	title={#2},
	boxed title size=title,
	boxed title style={%
			sharp corners,
			rounded corners=northwest,
			colback=tcbcolframe,
			boxrule=0pt,
		},
	underlay boxed title={%
			\path[fill=tcbcolframe] (title.south west)--(title.south east)
			to[out=0, in=180] ([xshift=5mm]title.east)--
			(title.center-|frame.east)
			[rounded corners=\kvtcb@arc] |-
			(frame.north) -| cycle;
		},
	#1
}{def}
\makeatother

%================================
% SOLUTION BOX
%================================

\makeatletter
\newtcolorbox{solution}{enhanced,
	breakable,
	colback=white,
	colframe=myg!80!black,
	attach boxed title to top left={yshift*=-\tcboxedtitleheight},
	title=Solution,
	boxed title size=title,
	boxed title style={%
			sharp corners,
			rounded corners=northwest,
			colback=tcbcolframe,
			boxrule=0pt,
		},
	underlay boxed title={%
			\path[fill=tcbcolframe] (title.south west)--(title.south east)
			to[out=0, in=180] ([xshift=5mm]title.east)--
			(title.center-|frame.east)
			[rounded corners=\kvtcb@arc] |-
			(frame.north) -| cycle;
		},
}
\makeatother

%================================
% Question BOX
%================================

\makeatletter
\newtcbtheorem{qstion}{Question}{enhanced,
	breakable,
	colback=white,
	colframe=mygr,
	attach boxed title to top left={yshift*=-\tcboxedtitleheight},
	fonttitle=\bfseries,
	title={#2},
	boxed title size=title,
	boxed title style={%
			sharp corners,
			rounded corners=northwest,
			colback=tcbcolframe,
			boxrule=0pt,
		},
	underlay boxed title={%
			\path[fill=tcbcolframe] (title.south west)--(title.south east)
			to[out=0, in=180] ([xshift=5mm]title.east)--
			(title.center-|frame.east)
			[rounded corners=\kvtcb@arc] |-
			(frame.north) -| cycle;
		},
	#1
}{def}
\makeatother

\newtcbtheorem[number within=chapter]{wconc}{Wrong Concept}{
	breakable,
	enhanced,
	colback=white,
	colframe=myr,
	arc=0pt,
	outer arc=0pt,
	fonttitle=\bfseries\sffamily\large,
	colbacktitle=myr,
	attach boxed title to top left={},
	boxed title style={
			enhanced,
			skin=enhancedfirst jigsaw,
			arc=3pt,
			bottom=0pt,
			interior style={fill=myr}
		},
	#1
}{def}



%================================
% NOTE BOX
%================================

\usetikzlibrary{arrows,calc,shadows.blur}
\tcbuselibrary{skins}
\newtcolorbox{note}[1][]{%
	enhanced jigsaw,
	colback=gray!20!white,%
	colframe=gray!80!black,
	size=small,
	boxrule=1pt,
	title=\textbf{Note:-},
	halign title=flush center,
	coltitle=black,
	breakable,
	drop shadow=black!50!white,
	attach boxed title to top left={xshift=1cm,yshift=-\tcboxedtitleheight/2,yshifttext=-\tcboxedtitleheight/2},
	minipage boxed title=1.5cm,
	boxed title style={%
			colback=white,
			size=fbox,
			boxrule=1pt,
			boxsep=2pt,
			underlay={%
					\coordinate (dotA) at ($(interior.west) + (-0.5pt,0)$);
					\coordinate (dotB) at ($(interior.east) + (0.5pt,0)$);
					\begin{scope}
						\clip (interior.north west) rectangle ([xshift=3ex]interior.east);
						\filldraw [white, blur shadow={shadow opacity=60, shadow yshift=-.75ex}, rounded corners=2pt] (interior.north west) rectangle (interior.south east);
					\end{scope}
					\begin{scope}[gray!80!black]
						\fill (dotA) circle (2pt);
						\fill (dotB) circle (2pt);
					\end{scope}
				},
		},
	#1,
}

%%%%%%%%%%%%%%%%%%%%%%%%%%%%%%
% SELF MADE COMMANDS
%%%%%%%%%%%%%%%%%%%%%%%%%%%%%%


\newcommand{\thm}[2]{\begin{Theorem}{#1}{}#2\end{Theorem}}
\newcommand{\cor}[2]{\begin{Corollary}{#1}{}#2\end{Corollary}}
\newcommand{\mlenma}[2]{\begin{Lenma}{#1}{}#2\end{Lenma}}
\newcommand{\mprop}[2]{\begin{Prop}{#1}{}#2\end{Prop}}
\newcommand{\clm}[3]{\begin{claim}{#1}{#2}#3\end{claim}}
\newcommand{\wc}[2]{\begin{wconc}{#1}{}\setlength{\parindent}{1cm}#2\end{wconc}}
\newcommand{\thmcon}[1]{\begin{Theoremcon}{#1}\end{Theoremcon}}
\newcommand{\ex}[2]{\begin{Example}{#1}{}#2\end{Example}}
\newcommand{\dfn}[2]{\begin{Definition}[colbacktitle=red!75!black]{#1}{}#2\end{Definition}}
\newcommand{\dfnc}[2]{\begin{definition}[colbacktitle=red!75!black]{#1}{}#2\end{definition}}
\newcommand{\qs}[2]{\begin{question}{#1}{}#2\end{question}}
\newcommand{\pf}[2]{\begin{myproof}[#1]#2\end{myproof}}
\newcommand{\nt}[1]{\begin{note}#1\end{note}}

\newcommand*\circled[1]{\tikz[baseline=(char.base)]{
		\node[shape=circle,draw,inner sep=1pt] (char) {#1};}}
\newcommand\getcurrentref[1]{%
	\ifnumequal{\value{#1}}{0}
	{??}
	{\the\value{#1}}%
}
\newcommand{\getCurrentSectionNumber}{\getcurrentref{section}}
\newenvironment{myproof}[1][\proofname]{%
	\proof[\bfseries #1: ]%
}{\endproof}

\newcommand{\mclm}[2]{\begin{myclaim}[#1]#2\end{myclaim}}
\newenvironment{myclaim}[1][\claimname]{\proof[\bfseries #1: ]}{}

\newcounter{mylabelcounter}

\makeatletter
\newcommand{\setword}[2]{%
	\phantomsection
	#1\def\@currentlabel{\unexpanded{#1}}\label{#2}%
}
\makeatother




\tikzset{
	symbol/.style={
			draw=none,
			every to/.append style={
					edge node={node [sloped, allow upside down, auto=false]{$#1$}}}
		}
}


% deliminators
\DeclarePairedDelimiter{\abs}{\lvert}{\rvert}
\DeclarePairedDelimiter{\norm}{\lVert}{\rVert}

\DeclarePairedDelimiter{\ceil}{\lceil}{\rceil}
\DeclarePairedDelimiter{\floor}{\lfloor}{\rfloor}
\DeclarePairedDelimiter{\round}{\lfloor}{\rceil}

\newsavebox\diffdbox
\newcommand{\slantedromand}{{\mathpalette\makesl{d}}}
\newcommand{\makesl}[2]{%
\begingroup
\sbox{\diffdbox}{$\mathsurround=0pt#1\mathrm{#2}$}%
\pdfsave
\pdfsetmatrix{1 0 0.2 1}%
\rlap{\usebox{\diffdbox}}%
\pdfrestore
\hskip\wd\diffdbox
\endgroup
}
\newcommand{\dd}[1][]{\ensuremath{\mathop{}\!\ifstrempty{#1}{%
\slantedromand\@ifnextchar^{\hspace{0.2ex}}{\hspace{0.1ex}}}%
{\slantedromand\hspace{0.2ex}^{#1}}}}
\ProvideDocumentCommand\dv{o m g}{%
  \ensuremath{%
    \IfValueTF{#3}{%
      \IfNoValueTF{#1}{%
        \frac{\dd #2}{\dd #3}%
      }{%
        \frac{\dd^{#1} #2}{\dd #3^{#1}}%
      }%
    }{%
      \IfNoValueTF{#1}{%
        \frac{\dd}{\dd #2}%
      }{%
        \frac{\dd^{#1}}{\dd #2^{#1}}%
      }%
    }%
  }%
}
\providecommand*{\pdv}[3][]{\frac{\partial^{#1}#2}{\partial#3^{#1}}}
%  - others
\DeclareMathOperator{\Lap}{\mathcal{L}}
\DeclareMathOperator{\Var}{Var} % varience
\DeclareMathOperator{\Cov}{Cov} % covarience
\DeclareMathOperator{\E}{E} % expected

% Since the amsthm package isn't loaded

% I prefer the slanted \leq
%\let\oldleq\leq % save them in case they're every wanted
%\let\oldgeq\geq
%\renewcommand{\leq}{\leqslant}
%\renewcommand{\geq}{\geqslant}

% % redefine matrix env to allow for alignment, use r as default
% \renewcommand*\env@matrix[1][r]{\hskip -\arraycolsep
%     \let\@ifnextchar\new@ifnextchar
%     \array{*\c@MaxMatrixCols #1}}


%\usepackage{framed}
%\usepackage{titletoc}
%\usepackage{etoolbox}
%\usepackage{lmodern}


%\patchcmd{\tableofcontents}{\contentsname}{\sffamily\contentsname}{}{}

%\renewenvironment{leftbar}
%{\def\FrameCommand{\hspace{6em}%
%		{\color{myyellow}\vrule width 2pt depth 6pt}\hspace{1em}}%
%	\MakeFramed{\parshape 1 0cm \dimexpr\textwidth-6em\relax\FrameRestore}\vskip2pt%
%}
%{\endMakeFramed}

%\titlecontents{chapter}
%[0em]{\vspace*{2\baselineskip}}
%{\parbox{4.5em}{%
%		\hfill\Huge\sffamily\bfseries\color{myred}\thecontentspage}%
%	\vspace*{-2.3\baselineskip}\leftbar\textsc{\small\chaptername~\thecontentslabel}\\\sffamily}
%{}{\endleftbar}
%\titlecontents{section}
%[8.4em]
%{\sffamily\contentslabel{3em}}{}{}
%{\hspace{0.5em}\nobreak\itshape\color{myred}\contentspage}
%\titlecontents{subsection}
%[8.4em]
%{\sffamily\contentslabel{3em}}{}{}  
%{\hspace{0.5em}\nobreak\itshape\color{myred}\contentspage}



%%%%%%%%%%%%%%%%%%%%%%%%%%%%%%%%%%%%%%%%%%%
% TABLE OF CONTENTS
%%%%%%%%%%%%%%%%%%%%%%%%%%%%%%%%%%%%%%%%%%%

\usepackage{tikz}
\definecolor{doc}{RGB}{0,60,110}
\usepackage{titletoc}
\contentsmargin{0cm}
\titlecontents{chapter}[3.7pc]
{\addvspace{30pt}%
	\begin{tikzpicture}[remember picture, overlay]%
		\draw[fill=doc!60,draw=doc!60] (-7,-.1) rectangle (-0.9,.5);%
		\pgftext[left,x=-3.5cm,y=0.2cm]{\color{white}\Large\sc\bfseries Chapter\ \thecontentslabel};%
	\end{tikzpicture}\color{doc!60}\large\sc\bfseries}%
{}
{}
{\;\titlerule\;\large\sc\bfseries Page \thecontentspage
	\begin{tikzpicture}[remember picture, overlay]
		\draw[fill=doc!60,draw=doc!60] (2pt,0) rectangle (4,0.1pt);
	\end{tikzpicture}}%
\titlecontents{section}[3.7pc]
{\addvspace{2pt}}
{\contentslabel[\thecontentslabel]{2pc}}
{}
{\hfill\small \thecontentspage}
[]
\titlecontents*{subsection}[3.7pc]
{\addvspace{-1pt}\small}
{}
{}
{\ --- \small\thecontentspage}
[ \textbullet\ ][]

\makeatletter
\renewcommand{\tableofcontents}{%
	\chapter*{%
	  \vspace*{-20\p@}%
	  \begin{tikzpicture}[remember picture, overlay]%
		  \pgftext[right,x=15cm,y=0.2cm]{\color{doc!60}\Huge\sc\bfseries \contentsname};%
		  \draw[fill=doc!60,draw=doc!60] (13,-.75) rectangle (20,1);%
		  \clip (13,-.75) rectangle (20,1);
		  \pgftext[right,x=15cm,y=0.2cm]{\color{white}\Huge\sc\bfseries \contentsname};%
	  \end{tikzpicture}}%
	\@starttoc{toc}}
\makeatother
%From M275 "Topology" at SJSU
\newcommand{\id}{\mathrm{id}}
\newcommand{\taking}[1]{\xrightarrow{#1}}
\newcommand{\inv}{^{-1}}

%From M170 "Introduction to Graph Theory" at SJSU
\DeclareMathOperator{\diam}{diam}
\DeclareMathOperator{\ord}{ord}
\newcommand{\defeq}{\overset{\mathrm{def}}{=}}

%From the USAMO .tex files
\newcommand{\ts}{\textsuperscript}
\newcommand{\dg}{^\circ}
\newcommand{\ii}{\item}

% % From Math 55 and Math 145 at Harvard
% \newenvironment{subproof}[1][Proof]{%
% \begin{proof}[#1] \renewcommand{\qedsymbol}{$\blacksquare$}}%
% {\end{proof}}

\newcommand{\liff}{\leftrightarrow}
\newcommand{\lthen}{\rightarrow}
\newcommand{\opname}{\operatorname}
\newcommand{\surjto}{\twoheadrightarrow}
\newcommand{\injto}{\hookrightarrow}
\newcommand{\On}{\mathrm{On}} % ordinals
\DeclareMathOperator{\img}{im} % Image
\DeclareMathOperator{\Img}{Im} % Image
\DeclareMathOperator{\coker}{coker} % Cokernel
\DeclareMathOperator{\Coker}{Coker} % Cokernel
\DeclareMathOperator{\Ker}{Ker} % Kernel
\DeclareMathOperator{\rank}{rank}
\DeclareMathOperator{\Spec}{Spec} % spectrum
\DeclareMathOperator{\Tr}{Tr} % trace
\DeclareMathOperator{\pr}{pr} % projection
\DeclareMathOperator{\ext}{ext} % extension
\DeclareMathOperator{\pred}{pred} % predecessor
\DeclareMathOperator{\dom}{dom} % domain
\DeclareMathOperator{\ran}{ran} % range
\DeclareMathOperator{\Hom}{Hom} % homomorphism
\DeclareMathOperator{\Mor}{Mor} % morphisms
\DeclareMathOperator{\End}{End} % endomorphism


\newcommand{\eps}{\epsilon}
\newcommand{\veps}{\varepsilon}
\newcommand{\ol}{\overline}
\newcommand{\ul}{\underline}
\newcommand{\wt}{\widetilde}
\newcommand{\wh}{\widehat}
\newcommand{\vocab}[1]{\textbf{\color{blue} #1}}
\providecommand{\half}{\frac{1}{2}}
\newcommand{\dang}{\measuredangle} %% Directed angle
\newcommand{\ray}[1]{\overrightarrow{#1}}
\newcommand{\seg}[1]{\overline{#1}}
\newcommand{\arc}[1]{\wideparen{#1}}
\DeclareMathOperator{\cis}{cis}
\DeclareMathOperator*{\lcm}{lcm}
\DeclareMathOperator*{\argmin}{arg min}
\DeclareMathOperator*{\argmax}{arg max}
\newcommand{\cycsum}{\sum_{\mathrm{cyc}}}
\newcommand{\symsum}{\sum_{\mathrm{sym}}}
\newcommand{\cycprod}{\prod_{\mathrm{cyc}}}
\newcommand{\symprod}{\prod_{\mathrm{sym}}}
\newcommand{\Qed}{\begin{flushright}\qed\end{flushright}}
\newcommand{\parinn}{\setlength{\parindent}{1cm}}
\newcommand{\parinf}{\setlength{\parindent}{0cm}}
% \newcommand{\norm}{\|\cdot\|}
\newcommand{\inorm}{\norm_{\infty}}
\newcommand{\opensets}{\{V_{\alpha}\}_{\alpha\in I}}
\newcommand{\oset}{V_{\alpha}}
\newcommand{\opset}[1]{V_{\alpha_{#1}}}
\newcommand{\lub}{\text{lub}}
\newcommand{\del}[2]{\frac{\partial #1}{\partial #2}}
\newcommand{\Del}[3]{\frac{\partial^{#1} #2}{\partial^{#1} #3}}
\newcommand{\deld}[2]{\dfrac{\partial #1}{\partial #2}}
\newcommand{\Deld}[3]{\dfrac{\partial^{#1} #2}{\partial^{#1} #3}}
\newcommand{\lm}{\lambda}
\newcommand{\uin}{\mathbin{\rotatebox[origin=c]{90}{$\in$}}}
\newcommand{\usubset}{\mathbin{\rotatebox[origin=c]{90}{$\subset$}}}
\newcommand{\lt}{\left}
\newcommand{\rt}{\right}
\newcommand{\bs}[1]{\boldsymbol{#1}}
\newcommand{\bd}{\mathbf}
\newcommand{\exs}{\exists}
\newcommand{\st}{\strut}
\newcommand{\dps}[1]{\displaystyle{#1}}

\newcommand{\sol}{\setlength{\parindent}{0cm}\textbf{\textit{Solution:}}\setlength{\parindent}{1cm} }
\newcommand{\solve}[1]{\setlength{\parindent}{0cm}\textbf{\textit{Solution: }}\setlength{\parindent}{1cm}#1 \Qed}
% Things Lie
\newcommand{\kb}{\mathfrak b}
\newcommand{\kg}{\mathfrak g}
\newcommand{\kh}{\mathfrak h}
\newcommand{\kn}{\mathfrak n}
\newcommand{\ku}{\mathfrak u}
\newcommand{\kz}{\mathfrak z}
\DeclareMathOperator{\Ext}{Ext} % Ext functor
\DeclareMathOperator{\Tor}{Tor} % Tor functor
\newcommand{\gl}{\opname{\mathfrak{gl}}} % frak gl group
\renewcommand{\sl}{\opname{\mathfrak{sl}}} % frak sl group chktex 6

% More script letters etc.
\newcommand{\SA}{\mathcal A}
\newcommand{\SB}{\mathcal B}
\newcommand{\SC}{\mathcal C}
\newcommand{\SF}{\mathcal F}
\newcommand{\SG}{\mathcal G}
\newcommand{\SH}{\mathcal H}
\newcommand{\OO}{\mathcal O}

\newcommand{\SCA}{\mathscr A}
\newcommand{\SCB}{\mathscr B}
\newcommand{\SCC}{\mathscr C}
\newcommand{\SCD}{\mathscr D}
\newcommand{\SCE}{\mathscr E}
\newcommand{\SCF}{\mathscr F}
\newcommand{\SCG}{\mathscr G}
\newcommand{\SCH}{\mathscr H}

% Mathfrak primes
\newcommand{\km}{\mathfrak m}
\newcommand{\kp}{\mathfrak p}
\newcommand{\kq}{\mathfrak q}

% number sets
\newcommand{\RR}[1][]{\ensuremath{\ifstrempty{#1}{\mathbb{R}}{\mathbb{R}^{#1}}}}
\newcommand{\NN}[1][]{\ensuremath{\ifstrempty{#1}{\mathbb{N}}{\mathbb{N}^{#1}}}}
\newcommand{\ZZ}[1][]{\ensuremath{\ifstrempty{#1}{\mathbb{Z}}{\mathbb{Z}^{#1}}}}
\newcommand{\QQ}[1][]{\ensuremath{\ifstrempty{#1}{\mathbb{Q}}{\mathbb{Q}^{#1}}}}
\newcommand{\CC}[1][]{\ensuremath{\ifstrempty{#1}{\mathbb{C}}{\mathbb{C}^{#1}}}}
\newcommand{\PP}[1][]{\ensuremath{\ifstrempty{#1}{\mathbb{P}}{\mathbb{P}^{#1}}}}
\newcommand{\HH}[1][]{\ensuremath{\ifstrempty{#1}{\mathbb{H}}{\mathbb{H}^{#1}}}}
\newcommand{\FF}[1][]{\ensuremath{\ifstrempty{#1}{\mathbb{F}}{\mathbb{F}^{#1}}}}
% expected value
\newcommand{\EE}{\ensuremath{\mathbb{E}}}
\newcommand{\charin}{\text{ char }}
\DeclareMathOperator{\sign}{sign}
\DeclareMathOperator{\Aut}{Aut}
\DeclareMathOperator{\Inn}{Inn}
\DeclareMathOperator{\Syl}{Syl}
\DeclareMathOperator{\Gal}{Gal}
\DeclareMathOperator{\GL}{GL} % General linear group
\DeclareMathOperator{\SL}{SL} % Special linear group

%---------------------------------------
% BlackBoard Math Fonts :-
%---------------------------------------

%Captital Letters
\newcommand{\bbA}{\mathbb{A}}	\newcommand{\bbB}{\mathbb{B}}
\newcommand{\bbC}{\mathbb{C}}	\newcommand{\bbD}{\mathbb{D}}
\newcommand{\bbE}{\mathbb{E}}	\newcommand{\bbF}{\mathbb{F}}
\newcommand{\bbG}{\mathbb{G}}	\newcommand{\bbH}{\mathbb{H}}
\newcommand{\bbI}{\mathbb{I}}	\newcommand{\bbJ}{\mathbb{J}}
\newcommand{\bbK}{\mathbb{K}}	\newcommand{\bbL}{\mathbb{L}}
\newcommand{\bbM}{\mathbb{M}}	\newcommand{\bbN}{\mathbb{N}}
\newcommand{\bbO}{\mathbb{O}}	\newcommand{\bbP}{\mathbb{P}}
\newcommand{\bbQ}{\mathbb{Q}}	\newcommand{\bbR}{\mathbb{R}}
\newcommand{\bbS}{\mathbb{S}}	\newcommand{\bbT}{\mathbb{T}}
\newcommand{\bbU}{\mathbb{U}}	\newcommand{\bbV}{\mathbb{V}}
\newcommand{\bbW}{\mathbb{W}}	\newcommand{\bbX}{\mathbb{X}}
\newcommand{\bbY}{\mathbb{Y}}	\newcommand{\bbZ}{\mathbb{Z}}

%---------------------------------------
% MathCal Fonts :-
%---------------------------------------

%Captital Letters
\newcommand{\mcA}{\mathcal{A}}	\newcommand{\mcB}{\mathcal{B}}
\newcommand{\mcC}{\mathcal{C}}	\newcommand{\mcD}{\mathcal{D}}
\newcommand{\mcE}{\mathcal{E}}	\newcommand{\mcF}{\mathcal{F}}
\newcommand{\mcG}{\mathcal{G}}	\newcommand{\mcH}{\mathcal{H}}
\newcommand{\mcI}{\mathcal{I}}	\newcommand{\mcJ}{\mathcal{J}}
\newcommand{\mcK}{\mathcal{K}}	\newcommand{\mcL}{\mathcal{L}}
\newcommand{\mcM}{\mathcal{M}}	\newcommand{\mcN}{\mathcal{N}}
\newcommand{\mcO}{\mathcal{O}}	\newcommand{\mcP}{\mathcal{P}}
\newcommand{\mcQ}{\mathcal{Q}}	\newcommand{\mcR}{\mathcal{R}}
\newcommand{\mcS}{\mathcal{S}}	\newcommand{\mcT}{\mathcal{T}}
\newcommand{\mcU}{\mathcal{U}}	\newcommand{\mcV}{\mathcal{V}}
\newcommand{\mcW}{\mathcal{W}}	\newcommand{\mcX}{\mathcal{X}}
\newcommand{\mcY}{\mathcal{Y}}	\newcommand{\mcZ}{\mathcal{Z}}


%---------------------------------------
% Bold Math Fonts :-
%---------------------------------------

%Captital Letters
\newcommand{\bmA}{\boldsymbol{A}}	\newcommand{\bmB}{\boldsymbol{B}}
\newcommand{\bmC}{\boldsymbol{C}}	\newcommand{\bmD}{\boldsymbol{D}}
\newcommand{\bmE}{\boldsymbol{E}}	\newcommand{\bmF}{\boldsymbol{F}}
\newcommand{\bmG}{\boldsymbol{G}}	\newcommand{\bmH}{\boldsymbol{H}}
\newcommand{\bmI}{\boldsymbol{I}}	\newcommand{\bmJ}{\boldsymbol{J}}
\newcommand{\bmK}{\boldsymbol{K}}	\newcommand{\bmL}{\boldsymbol{L}}
\newcommand{\bmM}{\boldsymbol{M}}	\newcommand{\bmN}{\boldsymbol{N}}
\newcommand{\bmO}{\boldsymbol{O}}	\newcommand{\bmP}{\boldsymbol{P}}
\newcommand{\bmQ}{\boldsymbol{Q}}	\newcommand{\bmR}{\boldsymbol{R}}
\newcommand{\bmS}{\boldsymbol{S}}	\newcommand{\bmT}{\boldsymbol{T}}
\newcommand{\bmU}{\boldsymbol{U}}	\newcommand{\bmV}{\boldsymbol{V}}
\newcommand{\bmW}{\boldsymbol{W}}	\newcommand{\bmX}{\boldsymbol{X}}
\newcommand{\bmY}{\boldsymbol{Y}}	\newcommand{\bmZ}{\boldsymbol{Z}}
%Small Letters
\newcommand{\bma}{\boldsymbol{a}}	\newcommand{\bmb}{\boldsymbol{b}}
\newcommand{\bmc}{\boldsymbol{c}}	\newcommand{\bmd}{\boldsymbol{d}}
\newcommand{\bme}{\boldsymbol{e}}	\newcommand{\bmf}{\boldsymbol{f}}
\newcommand{\bmg}{\boldsymbol{g}}	\newcommand{\bmh}{\boldsymbol{h}}
\newcommand{\bmi}{\boldsymbol{i}}	\newcommand{\bmj}{\boldsymbol{j}}
\newcommand{\bmk}{\boldsymbol{k}}	\newcommand{\bml}{\boldsymbol{l}}
\newcommand{\bmm}{\boldsymbol{m}}	\newcommand{\bmn}{\boldsymbol{n}}
\newcommand{\bmo}{\boldsymbol{o}}	\newcommand{\bmp}{\boldsymbol{p}}
\newcommand{\bmq}{\boldsymbol{q}}	\newcommand{\bmr}{\boldsymbol{r}}
\newcommand{\bms}{\boldsymbol{s}}	\newcommand{\bmt}{\boldsymbol{t}}
\newcommand{\bmu}{\boldsymbol{u}}	\newcommand{\bmv}{\boldsymbol{v}}
\newcommand{\bmw}{\boldsymbol{w}}	\newcommand{\bmx}{\boldsymbol{x}}
\newcommand{\bmy}{\boldsymbol{y}}	\newcommand{\bmz}{\boldsymbol{z}}

%---------------------------------------
% Scr Math Fonts :-
%---------------------------------------

\newcommand{\sA}{{\mathscr{A}}}   \newcommand{\sB}{{\mathscr{B}}}
\newcommand{\sC}{{\mathscr{C}}}   \newcommand{\sD}{{\mathscr{D}}}
\newcommand{\sE}{{\mathscr{E}}}   \newcommand{\sF}{{\mathscr{F}}}
\newcommand{\sG}{{\mathscr{G}}}   \newcommand{\sH}{{\mathscr{H}}}
\newcommand{\sI}{{\mathscr{I}}}   \newcommand{\sJ}{{\mathscr{J}}}
\newcommand{\sK}{{\mathscr{K}}}   \newcommand{\sL}{{\mathscr{L}}}
\newcommand{\sM}{{\mathscr{M}}}   \newcommand{\sN}{{\mathscr{N}}}
\newcommand{\sO}{{\mathscr{O}}}   \newcommand{\sP}{{\mathscr{P}}}
\newcommand{\sQ}{{\mathscr{Q}}}   \newcommand{\sR}{{\mathscr{R}}}
\newcommand{\sS}{{\mathscr{S}}}   \newcommand{\sT}{{\mathscr{T}}}
\newcommand{\sU}{{\mathscr{U}}}   \newcommand{\sV}{{\mathscr{V}}}
\newcommand{\sW}{{\mathscr{W}}}   \newcommand{\sX}{{\mathscr{X}}}
\newcommand{\sY}{{\mathscr{Y}}}   \newcommand{\sZ}{{\mathscr{Z}}}


%---------------------------------------
% Math Fraktur Font
%---------------------------------------

%Captital Letters
\newcommand{\mfA}{\mathfrak{A}}	\newcommand{\mfB}{\mathfrak{B}}
\newcommand{\mfC}{\mathfrak{C}}	\newcommand{\mfD}{\mathfrak{D}}
\newcommand{\mfE}{\mathfrak{E}}	\newcommand{\mfF}{\mathfrak{F}}
\newcommand{\mfG}{\mathfrak{G}}	\newcommand{\mfH}{\mathfrak{H}}
\newcommand{\mfI}{\mathfrak{I}}	\newcommand{\mfJ}{\mathfrak{J}}
\newcommand{\mfK}{\mathfrak{K}}	\newcommand{\mfL}{\mathfrak{L}}
\newcommand{\mfM}{\mathfrak{M}}	\newcommand{\mfN}{\mathfrak{N}}
\newcommand{\mfO}{\mathfrak{O}}	\newcommand{\mfP}{\mathfrak{P}}
\newcommand{\mfQ}{\mathfrak{Q}}	\newcommand{\mfR}{\mathfrak{R}}
\newcommand{\mfS}{\mathfrak{S}}	\newcommand{\mfT}{\mathfrak{T}}
\newcommand{\mfU}{\mathfrak{U}}	\newcommand{\mfV}{\mathfrak{V}}
\newcommand{\mfW}{\mathfrak{W}}	\newcommand{\mfX}{\mathfrak{X}}
\newcommand{\mfY}{\mathfrak{Y}}	\newcommand{\mfZ}{\mathfrak{Z}}
%Small Letters
\newcommand{\mfa}{\mathfrak{a}}	\newcommand{\mfb}{\mathfrak{b}}
\newcommand{\mfc}{\mathfrak{c}}	\newcommand{\mfd}{\mathfrak{d}}
\newcommand{\mfe}{\mathfrak{e}}	\newcommand{\mff}{\mathfrak{f}}
\newcommand{\mfg}{\mathfrak{g}}	\newcommand{\mfh}{\mathfrak{h}}
\newcommand{\mfi}{\mathfrak{i}}	\newcommand{\mfj}{\mathfrak{j}}
\newcommand{\mfk}{\mathfrak{k}}	\newcommand{\mfl}{\mathfrak{l}}
\newcommand{\mfm}{\mathfrak{m}}	\newcommand{\mfn}{\mathfrak{n}}
\newcommand{\mfo}{\mathfrak{o}}	\newcommand{\mfp}{\mathfrak{p}}
\newcommand{\mfq}{\mathfrak{q}}	\newcommand{\mfr}{\mathfrak{r}}
\newcommand{\mfs}{\mathfrak{s}}	\newcommand{\mft}{\mathfrak{t}}
\newcommand{\mfu}{\mathfrak{u}}	\newcommand{\mfv}{\mathfrak{v}}
\newcommand{\mfw}{\mathfrak{w}}	\newcommand{\mfx}{\mathfrak{x}}
\newcommand{\mfy}{\mathfrak{y}}	\newcommand{\mfz}{\mathfrak{z}}




\tikzsymbolsdefinesymbol{sus}{}{\begin{tikzpicture}[x=1.5,y=1.5] 
   \fill
   (5.44,4.97) .. controls (5.70,3.67) and (5.50,1.48) ..
   (5.36,0.86) .. controls (5.22,0.68) and (4.92,0.63) ..
   (4.24,0.58) .. controls (3.83,0.99) and (3.85,1.38) ..
   (3.83,1.69) .. controls (3.68,1.74) and (3.67,1.75) ..
   (3.46,1.69) .. controls (3.39,1.28) and (3.48,0.93) ..
   (3.42,0.52) .. controls (2.90,0.10) and (2.09,0.00) ..
   (1.80,0.21) .. controls (1.55,0.39) and (1.51,0.86) ..
   (1.54,1.56) .. controls (1.32,1.61) and (0.84,1.36) ..
   (0.52,1.83) .. controls (0.00,2.61) and (0.01,4.97) ..
   (0.16,5.23) .. controls (0.32,5.49) and (0.79,5.54) ..
   (1.09,5.60) .. controls (1.18,6.01) and (0.94,7.41) ..
   (2.37,7.75) .. controls (4.11,8.12) and (4.80,7.09) ..
   (4.89,6.71) .. controls (5.49,6.61) and (5.88,5.57) ..
   (5.44,4.97) -- cycle;
   \shade[top color=red, bottom color=red!40!black]
   (3.48,7.36) .. controls (3.76,7.33) and (4.10,7.04) ..
   (4.30,6.81) .. controls (3.35,6.76) and (2.20,6.73) ..
   (2.31,5.90) .. controls (2.36,5.54) and (2.44,3.40) ..
   (5.04,4.68) .. controls (5.35,2.16) and (5.01,2.05) ..
   (5.01,1.79) .. controls (5.01,1.61) and (5.14,1.48) ..
   (5.06,1.09) .. controls (4.83,1.01) and (4.57,0.88) ..
   (4.28,0.96) .. controls (4.05,1.35) and (4.28,1.51) ..
   (4.20,1.79) .. controls (4.41,1.84) and (4.54,1.92) ..
   (4.65,2.13) .. controls (4.02,1.97) and (3.19,2.16) ..
   (3.06,2.00) .. controls (2.93,1.77) and (3.09,1.38) ..
   (3.06,0.86) .. controls (2.72,0.47) and (2.42,0.44) ..
   (2.05,0.60) .. controls (1.89,1.22) and (1.68,2.62) ..
   (1.60,3.30) .. controls (1.27,5.48) and (1.47,6.39) ..
   (1.71,6.78) .. controls (1.97,7.43) and (3.19,7.38) ..
   (3.48,7.36) -- cycle
   (0.56,5.07) .. controls (0.77,5.17) and (0.93,5.15) ..
   (1.14,5.15) .. controls (1.03,4.34) and (1.08,2.78) ..
   (1.42,1.92) .. controls (1.14,1.84) and (1.06,1.95) ..
   (0.77,2.05) .. controls (0.43,2.88) and (0.43,4.58) ..
   (0.56,5.07) -- cycle;
   \shade[top color=blue!20, bottom color=blue!20!black]
   (2.84,6.11) .. controls (3.07,6.37) and (4.06,6.58) ..
   (4.61,6.37) .. controls (5.15,6.16) and (5.44,5.85) ..
   (5.21,5.38) .. controls (4.97,4.91) and (3.67,4.50) ..
   (3.15,4.94) .. controls (2.63,5.38) and (2.61,5.85) ..
   (2.84,6.11) -- cycle;
   \end{tikzpicture}}


%bullet size
\renewcommand\labelitemi{$\vcenter{\hbox{\tiny$\bullet$}}$}



\title{\Huge{Quantum Mechanics 1}\\PHYS	3143}
\author{\huge{Alex Feng}}
\date{}

\begin{document}
% Background color
%\pagestyle{empty}
\begin{titlepage}
\begin{tikzpicture}[remember picture,overlay]
\fill[PineGreen] (current page.south west) rectangle (current page.north east);


% Background Hexagon
\begin{scope}
\foreach \i in {2.5,...,22}
{\node[rounded corners,PineGreen!90,draw,regular polygon,regular polygon sides=6, minimum size=\i cm,ultra thick] at ($(current page.west)+(2.5,-5)$) {} ;}
\end{scope}

\foreach \i in {0.5,...,22}
{\node[rounded corners,PineGreen!90,draw,regular polygon,regular polygon sides=6, minimum size=\i cm,ultra thick] at ($(current page.north west)+(2.5,0)$) {} ;}

\foreach \i in {0.5,...,22}
{\node[rounded corners,PineGreen!98,draw,regular polygon,regular polygon sides=6, minimum size=\i cm,ultra thick] at ($(current page.north east)+(0,-9.5)$) {} ;}

\foreach \i in {21,...,6}
{\node[PineGreen!95,rounded corners,draw,regular polygon,regular polygon sides=6, minimum size=\i cm,ultra thick] at ($(current page.south east)+(-0.2,-0.45)$) {} ;}

% Title of the Report
\node[left,PineGreen!5,minimum width=0.625*\paperwidth,minimum height=3cm, rounded corners] at ($(current page.north east)+(0,-9.5)$){{\fontsize{25}{30} \selectfont \bfseries Quantum Mechanics 1}};

% Subtitle 
\node[left,PineGreen!10,minimum width=0.625*\paperwidth,minimum height=2cm, rounded corners] at ($(current page.north east)+(0,-11)$){{\huge PHYS 3143}};

% Author Name
\node[left,PineGreen!5,minimum width=0.625*\paperwidth,minimum height=2cm, rounded corners] at ($(current page.north east)+(0,-13)$){{\Large \textsc{Alex Feng}}};

% Year
\node[rounded corners,fill=PineGreen!95,text =PineGreen!5,regular polygon,regular polygon sides=6, minimum size=2.5 cm,inner sep=0,ultra thick] at ($(current page.west)+(2.5,-5)$) {\LARGE \bfseries 2024};

\end{tikzpicture}
\end{titlepage}

\newpage% or \cleardoublepage
% \pdfbookmark[<level>]{<title>}{<dest>}
\pdfbookmark[section]{\contentsname}{toc}
\tableofcontents
\pagebreak

\chapter{Stern-Gerlach Experiments}
\section{The Original Stern-Gerlach Experiment}

Particles have an intrintic spin $\mathbf{S}$, a form of angular momentum; although the spin angular momentum is real angular momentum, it is different from the classical orbital angular momentum. To understand the relationship between the angular momentum and magnetic moment of a charged particle, consider a classical of charge $q$ moving in a circular orbit of radius $\mathbf r$ and velocity $\mathbf v$. The magnetic moment (in Gaussian units)  $\bs{\mu}$ is given by
\[
	\boldsymbol{\mu} = \frac{I \mathbf A}{c} = \frac{q\mathbf r \times \mathbf b}{2c} = \frac{q}{2mc} \mathbf L
\] 
When it comes to intrinsic spin, we write
\[
	\bs{\mu} = \frac{gq}{2mc} \mathbf S
\] 
for a constant \(g\) experimentally determined for different particles. Changing the mass distribution of the classical model similarly results in such a factor, but it can not be applied to the intrinsic spin model. 

In the Stern-Gerlach experiment, a beam of silver atoms is directed through an inhomogeneous magnetic field. Recall that the force the particles experience is only dependent on field gradient and magnetic moment:
\[
\mathbf F = \nabla (\bs{\mu} \cdot \mathbf B) 
\] 
Surprisingly, the silver atoms only built up in two spots. This would mean that their magnetic moments could only take on two values (we now now that these correspond to the spins \(\pm \hbar/2\)).  

\section{Four Experiments}

\dfn{Ket vector}{
The symbol \(\ket{+\mathbf n}\) is a vector denoting positive spin in the positive \(\mathbf n\) direction. }
Suppose a particle exited an SG\textbf z device in state \(\ket{+ \mathbf z}\), i.e. \(S_z = + \hbar /2\). When a beam of these particles are sent through another SG\(\mathbf z\) device, all the particle will exit with the same state, \(\ket{+ \mathbf z}\).

Now, instead of sending the beam of particles in state \(\ket{+\mathbf z}\) through the SG\(\mathbf z\) device, send it through an SG\(\mathbf x\) device. We find that half the particles exit with state \(\ket{+ \mathbf x}\) and half with state \(\ket{- \mathbf x}\). Note that starting with particles in state \(\ket{+ \mathbf z}\) yields the same results. 

For the third experiment, send the particles in state \(\ket{+ \mathbf z}\) through an SG\(\mathbf x\) device, and send the resulting particles with state \(\ket{+ \mathbf x }\) through another SG\(\mathbf z\) device. Surprisingly, half the remaining particles have state \(\ket{+ \mathbf z}\) and the other half are in state \(\ket{- \mathbf z}\). This indicates that the Quantum world is different from our everyday classical experiences. 

Before moving on to the fourth experiment, consider the modified SG\(\mathbf n\) device, which sends paricles through one SG\(\mathbf n\) device, two SG\(\mathbf n\) devices with their polarities reversed, and another SG\(\mathbf n\) device at the end. Now, send particles in state \(\ket{+ \mathbf z}\) through the modified SG\(\mathbf x\) device, then directly through an SG\(\mathbf z\) device. Based on the previous experiments, we might expect half the particles to be in state \(\ket{+ \mathbf x}\); going through the SG\(\mathbf z\) device results in half of these \(\ket{+ \mathbf x}\) particles being in state \(\ket{+ \mathbf z}\). However, we find that all the exiting particles remained in state \(\ket{+ \mathbf z}\). We will later see that these effects can be explained using probability amplitudes. 

\section{The Quantum State Vector}

\dfn{Bra vector}{For all \(\ket{\psi}\), there exists a bra vector \(\bra{\psi},\)
which is an element of the dual space. The inner product of 
\(\bra{\psi}\) and \(\ket{\phi}\) is  \(\braket{\psi|\phi} \in \mathbb C \), which is the probability amplitude of finding a particle in state \(\ket{\phi}\) to be in state \(\ket{\psi}\).}
From the first experiment, we see that \(\braket{- \mathbf z | + \mathbf z} = 0\), so we can think of \(\ket{+ \mathbf z}\) and \(\ket{- \mathbf z}\) as orthogonal. Additionally, it will be convenient to require these state vectors to be unit vectors. An arbitrary spin state \(\ket{\psi}\) can be created by sending spin-\(\frac{1}{2}\) particles through an SG device oriented in an aritrary direction and selecting particles with positive spin. \(\ket{+ \mathbf z}\) and \(\ket{- \mathbf z}\) form a basis, thus we will describe the state \(\ket{\psi}\) as
\[
\ket{\psi} = c_+ \ket{+ \mathbf z} + c_- \ket{- \mathbf z}
\] 
Multiplying \(\bra{ \pm \mathbf z} \) shows that
\[
\ket{\psi} = \braket{+ \mathbf z | \psi} \ket{+ \mathbf z} + \braket{- \mathbf z | \psi} \ket{- \mathbf z}
\] 
Using the same technique, 
\[
\bra{\psi} = \braket{\psi |+\mathbf z } \bra{+ \mathbf z} + \braket{ \psi | - \mathbf z } \bra{- \mathbf z}
\] 
We will later find that \(\braket{\psi | \phi} = \braket{\phi | \psi}^*\). Using this, 
\[
\braket{\psi | \psi} = 1 = |\braket{+ \mathbf z| \psi}|^2 + |\braket{- \mathbf z | \psi}|^2
\] 
We interperet \(|\braket{+ \mathbf z | \psi}|^2\) to be the probability that the particle in state \(\ket \psi\) will be found in state \(\ket{ + \mathbf z}\). When \(\braket{+ \mathbf z | \psi}\) and \(\braket{- \mathbf z | \psi}\) are both non-zero, then the particle in state \(\ket \psi\) is actually in a superposition, instead of being in the definite state that classical mechanics expects. 

\section{Analysis of Experiment 3}

We know that \(|\braket{+ \mathbf z | + \mathbf x}|^2 =|\braket{- \mathbf z | + \mathbf x}|^2 = \frac{1}{2} \), thus, 
\[
	\braket{+ \mathbf z | + \mathbf x} = \frac{e^{i\delta_+}}{\sqrt 2} \text{ and } \braket{- \mathbf z | + \mathbf x} = \frac{e^{i\delta_-}}{\sqrt 2}
\]
We can use these values to calculate expected value \(\langle S_z \rangle\) and uncertainty \(\Delta S_z\):
\begin{align*}
	(\Delta S_z)^2 &= \langle ( S_z - \langle S_z \rangle ) ^2 \rangle\\
		       &=\langle S_z \rangle^2 - \langle S_z \rangle ^2
\end{align*}
\ex{Expected Value and Uncertainty}{Consider the spin-\(\frac{1}{2}\) particle in state
\[
\ket \psi = \frac{1}{2}\ket{+ \mathbf z} + \frac{i \sqrt 3}{2}\ket{- \mathbf z}
\] 
To find the expected value, 
\begin{align*}
\langle S_z \rangle & = |\braket{+ \mathbf z | \psi}|^2 \frac{\hbar}{2} - |\braket{- \mathbf z | \psi}|^2 \frac{\hbar}{2} \\
		    &=\frac{1}{4}\frac{\hbar}{2} - \frac{3}{4}\frac{\hbar}{2} \\
		    &= -\frac{\hbar}{4}
\end{align*}
To find the uncertainty, we first find \(\langle S_z^2 \rangle\),
\begin{align*}
\langle S_z \rangle & = |\braket{+ \mathbf z | \psi}|^2 \left(\frac{\hbar}{2}\right)^2 + |\braket{- \mathbf z | \psi}|^2 \left(-\frac{\hbar}{2}\right)^2 \\
		    &=\frac{1}{4}\frac{\hbar^2}{4} - \frac{3}{4}\frac{\hbar^2}{4} \\
		    &= \frac{\hbar^2}{4}
\end{align*}
It follows that 
\[
\Delta S_z = \frac{\sqrt 3}{4}\hbar
\] 

}
\section{Experiment 5}

	In this experiment, send particles with state \(\ket{+ \mathbf z}\) through an SG\(\mathbf x\) device, then make measurements of \(S_y\) on the particles exiting with state \(\ket{+ \mathbf x}\). From the third experiment, we already know that 50 percent of the particles have \(S_y = \hbar/2\) and 50 percent have \(S_y = - \hbar/2\). Another observer could be using different x/y/z axes and the results would be the same; this type of argument tells us that switching the SG\(\mathbf x\) and SG\(\mathbf y\) devices will not change the results of the third experiment. This argument has important implications. First, we already know that 
\[
	\ket{+ \mathbf y} = \frac{e^{i \gamma _+}}{\sqrt 2}\ket{+ \mathbf z} + \frac{e^{i \gamma _-}}{\sqrt 2}\ket{- \mathbf z} = \frac{e^{i \gamma_+}}{\sqrt 2}\left[ \ket{+ \mathbf z} + e^{i \gamma} \ket{ - \mathbf z}\right]
\] 
for \(\gamma = \gamma_- - \gamma_+\). Similarly, 
\[
	\ket{+ \mathbf x} = \frac{e^{i \delta_+}}{\sqrt 2} \left[ \ket{+ \mathbf z} + e^{i \delta} \ket{- \mathbf z} \right]
\] 
it follows that 
\[
\braket{+ \mathbf y | + \mathbf x} = \frac{e^{i(\delta_+ - \gamma_+)}}{2}\left[1 + e^{i(\delta - \gamma)}\right]
\] 
and finally,
\[
	|\braket{+ \mathbf y | + \mathbf x}|^2 = \frac{1}{2}\left[1 + \cos(\delta - \gamma) \right] 
\] 
\[
	|\braket{+ \mathbf y | + \mathbf x}|^2 = \frac{1}{2} \implies \delta - \gamma = \pm \frac{\pi}{2} 
\]
By convention, \(\delta = 0\), and if we ignore the phases \(\delta_+\) and \(\gamma_+\), we see that
\[
\ket{+ \mathbf x} = \frac{1}{\sqrt 2}\ket{+ \mathbf z} + \frac{1}{\sqrt 2}\ket{ - \mathbf z} 
\] 
and 
\[
\ket{+ \mathbf y} = \frac{1}{\sqrt 2} \ket{+ \mathbf z} + \frac{i}{\sqrt 2}\ket{ - \mathbf z}
\] 
Note that we have chosen \(\gamma = \pi/2\), which corresponds to a right-handed coordinate system. Results such as 
\[\ket{- \mathbf y}= \frac{1}{\sqrt2} \ket{+ \mathbf z} - \frac{i}{\sqrt 2} \ket{- \mathbf z}\]
arise from using \(\gamma = - \pi/2\) (left-handed coordinate system). 


\chapter{Rotation of Basis States and Matrix Mechanics}
\section{The Beginnings of Matrix Mechanics}
We can represent kets as column vectors in any basis we choose. For example, in the \(S_z\) basis, 
\[
\ket{\psi} \underset{S_z}\longrightarrow \begin{pmatrix} \braket{+ \mathbf z | \psi} \\ \braket{- \mathbf z | \psi} \end{pmatrix} 
\] 
Additionally, 
\[
\ket{+ \mathbf z} \underset{S_z}\longrightarrow \begin{pmatrix} 1 \\ 0 \end{pmatrix} \text{ and } \ket{- \mathbf z} \underset{S_z}\longrightarrow \begin{pmatrix} 0 \\ 1 \end{pmatrix} 
\] 
Similarly, we express bra vectors as row vectors:
\[
\bra{\psi} \underset{S_z}\longrightarrow  \begin{pmatrix} 

\braket{\psi | + \mathbf z} , \braket{\psi | -\mathbf z} 
\end{pmatrix} \] 
Thus, we can evaluate inner products using matrix mechanics, for example, 
\[
\braket{\psi | \psi} = \begin{pmatrix} \braket{\psi | + \mathbf z}, \braket{\psi | - \mathbf z} \end{pmatrix} \begin{pmatrix} \braket{+ \mathbf z | \psi} \\ \braket{ - \mathbf z | \psi} \end{pmatrix} = 1
\] 
\section{Rotation Operators}
\dfn{Rotation Operator}{The rotation operator \(\hat R (\theta \mathbf n)\) acts on kets by rotating them an angle \(\theta\) counterclockwise around axis \(\mathbf n\).}
Note that all operators will have a hat. Consider \(\hat R \left( \frac{\pi}{2} \mathbf j\right)\). This operator rotates kets 90 degrees clockwise around the y axis, so
\[
\ket{+ \mathbf x} = \hat R\left(\frac{\pi}{2}\mathbf j\right) \ket{+ \mathbf z}
\] 
We write the corresponding bra equation as
\[
\bra{+ \mathbf x} = \bra{+ \mathbf z}\hat R^ \dag \left(\frac{\pi}{2} \mathbf j\right)
\] 
where \(\hat R^\dag\) is the adjoint operator of \(\hat R\). Since 
\[
1 = \braket{+ \mathbf x | + \mathbf x} = \bra{+ \mathbf z} \hat R^\dag \left(\frac{\pi}{2} \mathbf j \right) \hat R \left(\frac{\pi}{2} \mathbf j \right) \ket{+ \mathbf z} = \braket{+ \mathbf z | + \mathbf z}
\]  
we know that \(\hat R^\dag\) and \(\hat R\) are inverses, and, in general, an operator \(\hat U\) satisfing \(U^\dag U = 1\) is called a unitary operator. 

\dfn{Hermitian Operator}{Any operator that is equal to its adjoint (self-adjoint) is called Hermitian.}
Now, consider an infinitesimal rotation operator, which we will express as
\[
\hat R(d\phi \mathbf k) = 1 - \frac{i}{\hbar}\hat J_z d\phi
\] 
First, notice that this form satisfies the conditions for rotation operators above. Also, note that \(\hat J_z\) must have the same dimentions as \(\hbar\), which are units of angular momentum. 
Taking the ajoint of \(\hat R(d\phi \mathbf k )\), we get
\[
\hat R^\dag(d \phi \mathbf k) = 1 + \frac{i}{\hbar}\hat J_z ^\dag d\phi
\] 
since we know that complex numbers are replaced by their conjugate when moving from kets to bras. We also know that the composition of the the rotation operator and its adjoint is 1. Thus
\begin{align*}
	\hat R^\dag (d \phi \mathbf k) \hat R (d \phi \mathbf k) &= \left(1 + \frac{i}{\hbar}\hat J_z^\dag d \phi \right) \left(1 - \frac{i}{\hbar} \hat J_z d\phi \right) \\
								 &=1 + \frac{i}{\hbar}\left(\hat J^\dag_z - \hat J_z \right) d \phi + O(d\phi ^2 ) = 1
\end{align*}
Thus, \(\hat J_z\) is Hermitian. To find the rotation operator for a finite angle \(\phi\), 
\[
d\phi = \lim_{N \to \infty}\frac{\phi}{N}
\] 
\[
	\hat R(\phi \mathbf k) = \lim_{N \to \infty}\left[ 1 - \frac{i}{\hbar}\hat J_z \left(\frac{\phi}{N}\right) \right]^N = e^{-i \hat J_z \phi/ \hbar}
\] 

Intuitively, \(\hat R(\phi \mathbf k) \ket{+ \mathbf z}\) does not change the state of \(\ket{+ \mathbf z}\), thus, \(\hat R(\phi \mathbf k)\) must only act as a overall phase factor. In general, if an operator acting on a state only changes the state by a constant factor, the state is called an eigenstate and the constant is the eigenvalue. Then, the Taylor series expansion of the rotation operator shows that \(\ket{+ \mathbf z}\) must be an eigenstate (or eigenket) of \(\hat J_z\). We claim that the eigenvalue of \(\ket{\pm \mathbf z}\) is \(\pm \hbar / 2\), which we will demonstrate is consistent later. Substituting \(\hat J_z\) with \(\hbar/2\) for \(\ket{+ \mathbf z}\) yields that 
\[
\hat R(\phi \mathbf k) \ket{+ \mathbf z} = e^{-i \phi /2} \ket{+ \mathbf z}
\] 
To show why the eigenvalues of \(\ket{+ \mathbf z}\) and \(\ket{- \mathbf z}\) are different for \(\hat J_z\), consider 
\[\ket{+ \mathbf x} = \frac{1}{\sqrt 2} \ket{+ \mathbf z} + \frac{1}{\sqrt 2} \ket{- \mathbf z}\]
\(\hat R (\phi \mathbf k)\) must not multiply \(\ket{+ \mathbf z}\) and \(\ket{- \mathbf z}\) by the same overall phase factor, otherwise \(\ket{+ \mathbf x}\) does not get rotated. Thus, \(\ket{+ \mathbf z}\) and \(\ket{- \mathbf z}\) must have different eigenvalues. Using 
\[
\hat J_z \ket{- \mathbf z} = - \frac{\hbar}{2} \ket{- \mathbf z}
\] 
we find that
\[
\hat R(\phi \mathbf k) \ket{- \mathbf z} = e^{i \phi / 2}\ket{- \mathbf z}
\] 
Now consider
\[
\hat R \left(\frac{\pi}{2} \mathbf k \right) \ket{+ \mathbf x} = e^{-i\pi/4} \left( \frac{1}{\sqrt 2} \ket{+ \mathbf z} + \frac{e^{i \pi/2}}{\sqrt 2} \ket{- \mathbf z} \right) = e^{-i \pi / 4} \ket{+ \mathbf y}
\] 
This result is correct, and only holds when \(\hat J_z \ket{\pm \mathbf z} = \pm \hbar / 2 \ket{\pm \mathbf z}\). Thus, the eigenvalues of the operator generating rotations about the z axis acting on \(\ket{\pm \mathbf z}\) are the values of the intrinsic spin angular momentum that characterizes these states! Finally, note that
\[
\hat R(2 \pi \mathbf k) \ket{+ \mathbf z} = e^{-i\pi} \ket{+ \mathbf z} = - \ket{+ \mathbf z} 
\]
and
\[
\hat R(2 \pi \mathbf k) \ket{- \mathbf z} = e^{i\pi} \ket{- \mathbf z} = - \ket{- \mathbf z} 
\]
Thus, rotating a spin-\(\frac{1}{2}\) state by 360 degrees and ending up where we started causes the state to pick up a minus sign. 


\section{The Identity and Projection Operators}
\dfn{Identity and Projection Operators}{The identity operator can be expressed as \[
1 = \ket{+ \mathbf z}\bra{+ \mathbf z} + \ket{- \mathbf z}\bra{- \mathbf z}
\] 
which are composed of the projection operators:
\[
\hat P_+ = \ket{+ \mathbf z}\bra{+ \mathbf z}\]
\[
\hat P _- = \ket{- \mathbf z}\bra{- \mathbf z}
\] 
}
We see that multiplying any spin-\(\frac{1}{2}\) state \(\ket{\psi}\) by the identity operator yields the original state \(\ket{\psi}\). Multiplying \(\ket{\psi}\) by the projection operators \(\hat P_+\) and \(\hat P_-\) yields the \(\ket{+ \mathbf z}\) and \(\ket{- \mathbf z}\) components of \(\ket{\psi}\) respectively. Note that the relation 
\[
\hat P_+ + \hat P_- = 1
\] 
is referred to as a completeness relation. We can think of the projection operators as a modified SG device that blocks certain pathways that a particle can take. 

It is convenient to introduce the following notation for basis states \(\ket{1}\) and \(\ket{2}\):
\[
\sum_i \ket{i}\bra{i} = 1 
\] 
for \(i = 1, 2\)

\section{Matrix Representations of Operators}
If we let 
\[
\hat A \ket{\psi} = \ket{\phi}
\] 
and we write 
\[
\ket{\psi} = \ket{+ \mathbf z} \braket{+ \mathbf z | \psi} + \ket{- \mathbf z} \braket{- \mathbf z | \psi}
\] 
\[
\ket{\phi} = \ket{+ \mathbf z} \braket{+ \mathbf z | \phi} + \ket{- \mathbf z}\braket{- \mathbf z | \phi}
\] 
then \(\hat A \ket{\psi} = \ket{\phi}\) becomes
\[
\hat A ( \ket{+ \mathbf z} \braket{+ \mathbf z | \psi} + \ket{- \mathbf z} \braket{- \mathbf z \psi}) = \ket{+ \mathbf z} \braket{+ \mathbf z | \phi} + \ket{- \mathbf z} \braket{- \mathbf z | \phi}
\] 
Multiplying by \(\bra{ \pm \mathbf z}\), we get
\[
\braket{+ \mathbf z | \hat A | + \mathbf z}\braket{+ \mathbf z | \psi} + \braket{+ \mathbf z | \hat A | - \mathbf z}\braket{- \mathbf z | \psi} = \braket{+ \mathbf z | \phi}
\] 
\[
\braket{- \mathbf z | \hat A | + \mathbf z}\braket{+ \mathbf z | \psi} + \braket{- \mathbf z | \hat A | - \mathbf z}\braket{- \mathbf z | \psi} = \braket{- \mathbf z | \phi}
\] 
We can express this in matrix form as
\[
	\begin{pmatrix} \braket{+ \mathbf z | \hat A | + \mathbf z} & \braket{+ \mathbf z| \hat A| - \mathbf z} \\ \braket{- \mathbf z| \hat A| + \mathbf z} & \braket{- \mathbf z | \hat A | - \mathbf z} \end{pmatrix} \begin{pmatrix} \braket{+ \mathbf z | \psi} \\ \braket{- \mathbf z | \psi} \end{pmatrix} = \begin{pmatrix} \braket{+ \mathbf z | \phi} \\ \braket{- \mathbf z | \phi} \end{pmatrix} 
\] 
which is represented in the \(S_z\) basis. We indicate a representation through an arrow:
\[
	\hat A \underset{S_z}\longrightarrow \begin{pmatrix} \braket{+ \mathbf z| \hat A | + \mathbf z} & \braket{+\mathbf z |\hat A| - \mathbf z}\\ \braket{- \mathbf z| \hat A| + \mathbf z} & \braket{- \mathbf z | \hat A| - \mathbf z} \end{pmatrix} = \begin{pmatrix} A_{11} & A_{12} \\ A_{21} & A_{22} \end{pmatrix} 
\] 
In general, for basis states \(\ket{1}\) and \(\ket{2}\) we can express matrix elements \(A_{ij}\) as
\[
A_{ij} = \braket{i | \hat A |j}
\] 

We can express \(\hat P_+\) as
\[
	\hat P_+ \underset{S_z}\longrightarrow \begin{pmatrix} 1 & 0 \\ 0 & 0 \end{pmatrix} 
\] 
and similarly, 
\[
	\hat P_- \underset{S_z}\longrightarrow \begin{pmatrix} 0 & 0 \\ 0 & 1 \end{pmatrix} 
\] 
Thus, the completeness relation in matrix form becomes
\[
	\begin{pmatrix} 1 & 0 \\ 0 & 0 \end{pmatrix} + \begin{pmatrix} 0 & 0 \\ 0 & 1 \end{pmatrix} = \mathbb I
\] 
We can also express \(\hat J_z\) in matrix form by using its eigenvalues:
\[
	\hat J_z \underset{S_z}\longrightarrow \begin{pmatrix} \braket{+ \mathbf z | \hat J_z | + \mathbf z} & \braket{+ \mathbf z | \hat J_z | - \mathbf z } \\ \braket{- \mathbf z | \hat J_z | + \mathbf z} & \braket{- \mathbf z | \hat J_z | - \mathbf z} \end{pmatrix} = \begin{pmatrix} \hbar / 2 & 0 \\ 0 & -\hbar / 2 \end{pmatrix} 
\] 

To find the form of the matrix representing an adjoint, consider the operator \(\hat A\), and states \(\chi\), \(\psi\), and \(\phi\). 
\[
\braket{\chi | \hat A | \psi} = \braket{\chi | \phi}
\] 
\[
\braket{\psi | \hat A^ \dag | \chi} = \braket{\phi | \chi}
\] 
Thus, we find that 
\[
\braket{\psi | \hat A^ \dag | \chi} = \braket{\chi | \hat A | \psi}^*
\] 
It follows that 
\[
\braket{i | \hat A ^\dag | j} = \braket{ j | \hat A | i}^*
\] 
\dfn{Adjoint Matrix}{We define the adjoint matrix \(\mathbb A^\dag\) as the transpose conjugate of the matrix \(\mathbb A\), i.e., \(A^\dag_{ij} = A^*_{ji}.\)}
Since a the matrix representation of a Hermitian operator equals its conjugate transpose by definition, operators such as \(\hat P_+\) or \(\hat J_z\) must be Hermitian.

To find the product of two operators \(\hat A\) and \(\hat B\), we first form the matrix element
\[
\braket{i | \hat A \hat B | j}
\] 
Inserting the identity element, we find that
\[
\braket{i | \hat A \hat B | j} = \bra i \hat A \left( \sum_k \ket{k} \bra{k} \right) \hat B \ket{j} = \sum_k \braket{i | \hat A | k}\braket{k | \hat B | j} \]
which is the same as usual matrix multiplication. Intuitively, 
\[
	(\hat A \hat B)^\dag = \hat B^\dag \hat A^\dag
\] 
\section{Changing Representations}
Consider the following equation:
\[
	\ket{\psi '} \underset{S_z}\longrightarrow \begin{pmatrix} \braket{+ \mathbf z | \psi '} \\ \braket{- \mathbf z | \psi '} \end{pmatrix} = \begin{pmatrix} \braket{+ \mathbf z | \hat R^ \dag ( \frac{\pi}{2} \mathbf j ) | \psi } \\ \braket{- \mathbf z | \hat R^ \dag ( \frac{\pi}{2} \mathbf j ) |\psi } \end{pmatrix} = \begin{pmatrix} \braket{+ \mathbf x | \psi} \\ \braket{- \mathbf x | \psi} \end{pmatrix} \underset{S_x}\longleftarrow \ket{\psi} 
\] 
The subtlety is that left side involes the state \(\ket{\psi}\) being rotated (active transformation); the right side simply has a rotated basis (passive transformation). 

In general, we can change the basis of any state by using the rotation operator \(\hat R ^\dag\), for example, 
\begin{align*}
	\begin{pmatrix} \braket{+ \mathbf x | \psi} \\ \braket{- \mathbf x | \psi} \end{pmatrix} &= \begin{pmatrix} \braket{+ \mathbf x | + \mathbf z} & \braket{+ \mathbf x | - \mathbf z}\\ \braket{- \mathbf x | + \mathbf z} & \braket{- \mathbf x | - \mathbf z} \end{pmatrix} \begin{pmatrix} \braket{+ \mathbf z | \psi} \\ \braket{- \mathbf z |\psi} \end{pmatrix} \\
												 &=\begin{pmatrix} \braket{+ \mathbf z | \hat R^\dag ( \frac{\pi}{2} \mathbf j )| + \mathbf z} & \braket{+ \mathbf z | \hat R^\dag ( \frac{\pi}{2} \mathbf j )| - \mathbf z} \\ \braket{- \mathbf z | \hat R^\dag ( \frac{\pi}{2} \mathbf j )| + \mathbf z}  & \braket{- \mathbf z | \hat R^\dag ( \frac{\pi}{2} \mathbf j )| - \mathbf z}  \end{pmatrix} \begin{pmatrix} \braket{+ \mathbf z | \psi} \\ \braket{- \mathbf z | \psi} \end{pmatrix} 
\end{align*}
We call the matrix in the example above \(\mathbb S^\dag\) (or more specifically \(\mathbb S^\dag (\frac{\pi}{2} \mathbf j)\)). Let \(\hat A\) be an operator represented as a matrix in the \(S_z\) basis. We can transform it to the \(S_x \) basis by using the \(\mathbb S^\dag\) from above and writing
\[
\hat A \underset{S_x} \longrightarrow \mathbb S^\dag \mathbb A \mathbb S
\] 
%insert example here later
\section{Expectation Values}
Observe that
\begin{align*}
	\langle S_z \rangle &= \frac{\hbar}{2} |\braket{+ \mathbf z | \psi }|^2 - \frac{\hbar}{2} |\braket{- \mathbf z | \psi}|^2\\
			    &=\left( \braket{\psi | + \mathbf z}, \braket{\psi | - \mathbf z} \right) \begin{pmatrix} \frac{\hbar}{2} & 0 \\ 0 & -\frac{\hbar}{2}  \end{pmatrix} \begin{pmatrix} \braket{+ \mathbf z | \psi} \\ \braket{- \mathbf z | \psi} \end{pmatrix}  
\end{align*}
Thus, we can simply express expected value in the form
\[
\langle S_z \rangle = \braket{\psi | \hat J_z | \psi}
\]
\chapter{Angular Momentum}
\section{Rotations Do Not Commute and Neither Do the Generators}
Consider the matrix equation
\[
	\begin{pmatrix} A'_x \\ A'_y \\ A'_z \end{pmatrix} = \begin{pmatrix} \cos \phi & - \sin \phi & 0 \\ \sin \phi & \cos \phi & 0 \\ 0& 0 & 1 \end{pmatrix}  \begin{pmatrix} A_x \\ A_y \\ A_z \end{pmatrix} 
\] 
which rotates the vector by an angle \(\phi\) counterclockwise about the z axis; we express this matrix as 
\[
\mathbb S (\phi \mathbf k) =  \begin{pmatrix} \cos \phi & - \sin \phi & 0 \\ \sin \phi & \cos \phi & 0 \\ 0& 0 & 1 \end{pmatrix}  \
\] 
For a small angle \(\Delta \phi\), we approximate this as
\[
	\mathbb S (\Delta \phi \mathbf k) = \begin{pmatrix} 1 - \Delta \phi ^2 / 2 & - \Delta \phi & 0 \\ \Delta \phi & 1 - \Delta \phi ^2 /2 & 0 \\ 0 & 0 & 1 \end{pmatrix}  
\]
rotations about other axes can be obtained from the cyclic substitution \(x \to y, y \to z, z \to y\). Thus, 
\[
	\mathbb S(\Delta \phi \mathbf i) = \begin{pmatrix} 1 & 0 & 0 \\ 0 & \cos \Delta \phi & - \sin \Delta \phi \\ 0 & \sin \Delta \phi & \cos \Delta \phi \end{pmatrix} \approx \begin{pmatrix} 1 & 0 & 0 \\ 0 & 1- \Delta \phi^2 /2 & - \Delta \phi \\ 0 & \Delta \phi & 1 - \Delta \phi ^2 /2 \end{pmatrix} 
\] 
Similarly, 
\[
	\mathbb S ( \Delta \phi \mathbf j) = \begin{pmatrix} 1- \Delta \phi^2/2 & 0 & \Delta \phi \\ 0 & 1 & 0 \\ - \Delta \phi & 0 & 1 - \Delta \phi ^2 /2 \end{pmatrix} 
\] 
We know that the \(\mathbb S\) matrices are representation of the rotation operators. Thus
\[
\mathbb S ( \Delta \phi \mathbf i) \mathbb S (\Delta \phi \mathbf j) -\mathbb S ( \Delta \phi \mathbf j) \mathbb S (\Delta \phi \mathbf i) = \begin{pmatrix} 0 & - \Delta \phi ^2 & 0 \\ \Delta \phi ^2 & 0 & 0 \\ 0 & 0 & 0 \end{pmatrix} = \mathbb S (\Delta \phi ^2 \mathbf k) - \mathbb I
\]
It follows that 
\begin{align*}
&\mathrel{\phantom{=}} \hat R(\Delta \phi \mathbf i) \hat R(\Delta \phi \mathbf j) - \hat R(\Delta \phi \mathbf j)\hat R(\Delta \phi \mathbf i)\\
&=e^{i \hat J _x \phi / \hbar} e^{-i \hat J _y \phi / \hbar} - e^{i \hat J_y \phi / \hbar} e^{-i \hat J _x \phi / \hbar} \\
\begin{split}
&\approx \left\{ 1 - \frac{i \hat J_x \Delta \phi}{\hbar} - \frac{1}{2}\left(\frac{\hat J_x \Delta \phi}{\hbar}\right)^2\right\} \left\{ 1 - \frac{i \hat J_y \Delta \phi}{\hbar} - \frac{1}{2}\left(\frac{\hat J_y \Delta \phi}{\hbar}\right)^2\right\}\\
& \qquad - \left\{ 1 - \frac{i \hat J_y \Delta \phi}{\hbar} - \frac{1}{2}\left(\frac{\hat J_y \Delta \phi}{\hbar}\right)^2\right\} \left\{ 1 - \frac{i \hat J_x \Delta \phi}{\hbar} - \frac{1}{2}\left(\frac{\hat J_x \Delta \phi}{\hbar}\right)^2\right\} 
\end{split}\\
&\approx  \frac{(\hat J_y \hat J_x - \hat J_x \hat J_y)\Delta \phi^2 }{\hbar^2}\\
&= \left( 1 - \frac{i \hat J_z \Delta \phi ^2}{\hbar} \right) -1
\end{align*}
\dfn{Commutator}{We define the commutator of two operators \(\hat A\) and \(\hat B\) to be 
\[
 [ \hat A, \hat B ] = \hat A \hat B - \hat B \hat A
\]}
Thus, we find that \[
\hat J_x  \hat J_y - \hat J_y \hat J_x = i \hbar \hat J_z \implies [\hat J_x , \hat J_y ] = i \hbar \hat J_z
\] 
By cyclic substitution, 
\[
[\hat J_y, \hat J_z ] = i \hbar \hat J_x
\] 
\[
[\hat J_z, \hat J_x] = i \hbar \hat J_y
\] 
These important commutation relations are general and also apply for orbital angular momentum operators. 
%\section{Commuting Operators}


\end{document}
